\chapter{财报分析与数据提取}
\chapterlead{对应主选题“财报分析”与子选题“数据提取”。一张真值表是整个系统的地基，这一章讲它是怎么建起来的。}

\section{从 PDF 到真值表}

年报 PDF 的版式五花八门。同一个“营业收入”，可能出现在“主要会计数据和财务指标”表、“利润表及现金流量表相关科目变动分析表”和合并利润表里，
数字相同但口径说明不同；合并报表与母公司报表的科目名称几乎一样；一张资产负债表跨两三页，续页没有表头。
我们没有试图写一个“万能解析器”，而是让两条出错方式完全不同的通道各自取数，再互相核对（图~\ref{fig:fin-dual}）。

\begin{figure}[htbp]
\centering
\begin{tikzpicture}[node distance=4mm and 6mm]
  \node[doc,minimum width=1.9cm,minimum height=1.1cm,font=\small] (pdf) {年报 PDF};
  \node[prog,right=8mm of pdf,yshift=8mm,minimum width=2.6cm] (loc) {\code{locate}\\ \scriptsize 关键页、单位、扫描页};
  \node[agent,right=of loc,minimum width=2.6cm] (llm) {分析师会话\\ \scriptsize 按页读，写 facts.json};
  \node[prog,right=of llm,minimum width=2.5cm] (ev) {证据校验\\ \scriptsize 所写页码逐字查找};
  \node[prog,right=8mm of pdf,yshift=-10mm,minimum width=2.6cm] (parse) {表格解析\\ \scriptsize 稀疏列 · 续表 · 折行};
  \node[prog,right=of parse,minimum width=2.6cm] (rule) {规则通道\\ \scriptsize 指标别名 + 表头期间};
  \node[gate,right=11mm of rule,yshift=9mm] (m) {合并};
  \node[prog,right=8mm of m,fill=progc!18,minimum width=1.6cm] (tt) {\textbf{真值表}};
  \node[gate,below=6mm of tt] (tie) {勾稽};
  \node[doc,below=6mm of m,text width=2.2cm,font=\scriptsize] (iss) {冲突 / 拒收 / 勾稽失败\\ → 阻断估值，退回核对};
  \draw[arr] (pdf) |- (loc); \draw[arr] (loc) -- (llm); \draw[arr] (llm) -- (ev);
  \draw[arr] (pdf) |- (parse); \draw[arr] (parse) -- (rule);
  \draw[arr] (ev) -| (m); \draw[arr] (rule) -| (m);
  \draw[arr] (m) -- (tt); \draw[arr] (tt) -- (tie);
  \draw[darr] (m) -- (iss); \draw[darr] (tie) -| (iss.south);
\end{tikzpicture}
\caption{双通道抽取。模型通道与规则通道独立取数，合并时一致则互为佐证，不一致则记为问题并阻断下游估值。}
\label{fig:fin-dual}
\end{figure}

\subsection{模型通道：抄数必须带页码}

分析师会话先运行 \code{locate} 拿到关键页清单（合并报表各在哪几页、母公司报表在哪几页、金额单位是元还是万元、有没有扫描页），
然后按页阅读，把指标写进 \file{facts.json}。每条记录包括指标编号、期间、原文数字串、单位、页码、原文科目标签和口径说明。
指标字典共 50 个指标，覆盖利润表、资产负债表、现金流量表、主要会计数据与非经常性损益，另有 8 个由程序派生的指标（如毛利率、净现比），
派生指标不允许模型填写。

\code{verify-facts} 对每一条记录做证据校验：到所写页码的原文里，按数字边界查找原文数字串。
“按数字边界”的意思是，抄“1,688.38”时原文里的“11,688.38”或“1,688.385”都不算找到。找不到就拒收，并告诉模型是哪一条、在哪一页没找到。
分析师会话据此改正页码或数字后重跑，直到拒收为 0。这一步的价值在实际运行中很明显：宇树科技的分析师会话在迭代中先后 6 次收到“未在所写页码的原文中找到该数字”，
每次都是页码写错或数字抄错，改正后才通过。

\subsection{规则通道：对付真实年报的版式}

规则通道不调用模型，用 pdfplumber 解析表格，按指标别名匹配科目行、按表头识别期间列。在虚构的样例年报上它一直工作得很好，
但第一次跑真实年报时，三家公司的 301 个模型核验值中，规则通道只独立取到并核对一致了 115 个，宁德时代一个也没取到。
逐页排查后，我们找到了三类版式问题，分别写了处理逻辑。

\begin{description}[leftmargin=0pt,labelsep=0.5em,font=\heiti\bfseries\color{ink}]
  \item[稀疏网格] 有些年报的表格线条稀疏，解析器把一个逻辑列切成几列：科目在第 2 列，表头“期末余额”比下面的数字向右错开一列，其余列全空。宁德时代的报表几乎全是这种版式。处理办法是合并全空列，并把表头与数字按位置重新对齐。
  \item[跨页续表] 合并报表常跨两三页，续页没有表头，规则通道找不到期间列就整页跳过。茅台的营业利润、存货、资产总计、经营现金流都在续页上。处理办法是识别“本页第一张表、上方无标题、上一页末尾无新标题”的续表，把上一页的表头接过来；两页的列数对不上时，宁可不认作续表。
  \item[科目折行] 长科目名会折成两三行，数字落在第一行，例如“归属于上市公司股东的扣除非经常性损益的净利润”。处理办法是把没有数字的后续行接回上一行。
\end{description}

三项处理的效果见表~\ref{tab:fin-ablation} 与图~\ref{fig:fin-ablation}。我们逐项关闭再重跑，发现它们的作用强烈依赖版式：
去掉稀疏列合并，宁德时代的佐证数从 101 掉到 0；茅台和招商轮船的增量几乎全部来自续表链接；折行合并只对宁德时代有效。
更重要的是，七种配置下两条通道同时给出值、但值不同的情况一次也没有出现。每一项处理都只让规则通道多取到，没有让它取错。

\begin{table}[htbp]
\centering
\caption{规则通道版式处理的逐项消融（佐证数，即规则通道独立取到且与模型一致的数值个数）}
\label{tab:fin-ablation}
\small
\begin{tabular}{@{}lrrrrrrr@{}}
\toprule
 & \multicolumn{4}{c}{\hd{逐项加上}} & \multicolumn{3}{c}{\hd{从完整配置去掉一项}} \\
\cmidrule(lr){2-5}\cmidrule(l){6-8}
\hd{公司（模型核验值）} & 全关 & +稀疏列 & +续表 & 完整 & $-$稀疏列 & $-$续表 & $-$折行 \\
\midrule
贵州茅台（98） & 60 & 60 & 90 & 90 & 86 & 60 & 90 \\
\rowband
宁德时代（106） & 0 & 25 & 69 & 101 & 0 & 47 & 69 \\
招商轮船（97） & 55 & 55 & 89 & 89 & 89 & 55 & 89 \\
\midrule
\rowband
合计（301） & 115 & 140 & 248 & \textbf{280} & 175 & 162 & 248 \\
佐证率 & 38.2\% & 46.5\% & 82.4\% & \textbf{93.0\%} & 58.1\% & 53.8\% & 82.4\% \\
\bottomrule
\end{tabular}
\end{table}

\begin{figure}[htbp]
\centering
\begin{tikzpicture}
\begin{axis}[ybar, bar width=8pt, width=0.95\linewidth, height=5.0cm, ymin=0, ymax=118, ylabel={佐证数},
  symbolic x coords={a,b,c,d,e,f,g}, xtick=data,
  xticklabels={全关,+稀疏列,+续表,完整,$-$稀疏列,$-$续表,$-$折行},
  legend style={at={(0.01,0.98)},anchor=north west}, legend columns=3,
  enlarge x limits=0.08, ymajorgrids, grid style={rule!60}, axis x line*=bottom, axis y line*=left]
\addplot[fill=ink!80,draw=ink] coordinates {(a,60)(b,60)(c,90)(d,90)(e,86)(f,60)(g,90)};
\addplot[fill=seal!75,draw=seal] coordinates {(a,0)(b,25)(c,69)(d,101)(e,0)(f,47)(g,69)};
\addplot[fill=gold!70,draw=gold] coordinates {(a,55)(b,55)(c,89)(d,89)(e,89)(f,55)(g,89)};
\legend{贵州茅台,宁德时代,招商轮船}
\end{axis}
\end{tikzpicture}
\caption{版式处理的作用因公司而异。七种配置下两通道错配均为 0。}
\label{fig:fin-ablation}
\end{figure}

\subsection{两通道合并}

合并时，同一指标、同一期间，两条通道的值一致（容差按原文展示精度）就记为 \code{rule+llm}，置信度高，所有出处都保留；
只有一方有值就记录该方；两方都有值但不一致，记为问题，阻断估值。三家年报的最终结果见表~\ref{tab:fin-real}：
301 个模型核验值全部通过证据校验，280 个得到规则通道独立佐证，冲突 0，56 项勾稽全部通过。

\begin{table}[htbp]
\centering
\caption{三家真实年报的双通道结果（当前代码重跑）}
\label{tab:fin-real}
\small
\begin{tabular}{@{}lrrrrrrrr@{}}
\toprule
\hd{公司} & \hd{模型条目} & \hd{拒收} & \hd{规则事实} & \hd{佐证} & \hd{仅模型} & \hd{仅规则} & \hd{冲突} & \hd{勾稽} \\
\midrule
贵州茅台 & 98 & 0 & 136 & 90 & 8 & 18 & 0 & 20/20 \\
\rowband
宁德时代 & 106 & 0 & 119 & 101 & 5 & 8 & 0 & 15/15 \\
招商轮船 & 97 & 0 & 144 & 89 & 8 & 23 & 0 & 21/21 \\
\midrule
合计 & 301 & 0 & 399 & 280 & 21 & 49 & 0 & 56/56 \\
\bottomrule
\end{tabular}
\end{table}

“仅模型”的 21 个值集中在规则通道已知不覆盖的地方，例如股本、部分扣非季度别名；“仅规则”多为上上年数据和财报列示的同比。
需要说明的是，“零冲突”不等于“抽取完全正确”，两条通道理论上可能一致地错。但一个是表格结构错误，一个是阅读错误，同一格上犯同一个错的概率很低，
再加上勾稽的约束，我们认为真值表在这三家公司上是可靠的。三家年报只是开发集，更多公司的评测正在进行（第 13 章）。

\subsection{经营数据与季报}

年报里的分产品、分地区、产销量、产能等经营数据表，由分析师会话整表抄进 \file{extras.json}，\code{verify-extra} 逐格回原文核对，派生的占比、单价、毛利率由程序计算。
五个工作区的经营数据各有 17～20 张表、126～367 个单元格，最终拒收均为 0。

季报通过 \code{--supplement} 作为补充文件接入，只走模型通道。原因是规则通道还不能识别季报的表头（“本报告期”“年初至报告期末”会被当成年度），
这是已知问题，列在第 13 章的待做事项里。

\section{真值表上的分析}

真值表是一个“指标 × 期间”的表，每格带数值、单位、出处和取得方式。财报分析的全部计算都在这张表上进行，不再回头读 PDF。

\subsection{同比、环比与单季拆分}

同比和环比由程序计算，存为百分数值。A 股定期报告只披露累计数，研究员最关心的“第四季度发生了什么”在任何一份报告里都不完整存在，
需要用全年数减前三季度累计数得到。系统按五条规则（$\text{Q2}=\text{H1}-\text{Q1}$、$\text{Q3}=\text{Q1\text{-}3}-\text{H1}$、$\text{Q4}=\text{全年}-\text{Q1\text{-}3}$ 等）反复推导，直到没有新的单元格可以算出。
规则只作用于发生额：资产负债表科目是时点数，每股收益的分母各期不同，比率不可加，这几类都不拆分。

\subsection{勾稽校验}

勾稽分四组（表~\ref{tab:fin-tie}）。其中“列示同比”这一组几乎不费力，却对最难发现的两类抽取错误特别敏感：本期列与上期列看反时同比的符号会反过来，
单位错一个量级时同比会差得离谱。

\begin{table}[htbp]
\centering
\caption{勾稽检查}
\label{tab:fin-tie}
\small
\begin{tabularx}{\linewidth}{@{}l L L@{}}
\toprule
\hd{组} & \hd{核对内容} & \hd{能发现的抽取错误} \\
\midrule
资产负债表 & 资产总计 = 负债和所有者权益总计；所有者权益 = 归母权益 + 少数股东权益 & 母公司与合并混取、续表错列、单位错一个量级 \\
\rowband
非经常性损益 & 列示合计 = 归母净利润 − 扣非归母净利润 & 合计取错表、扣非取错口径 \\
列示同比 & 财报“本年比上年增减”与程序计算的同比一致 & 本期与上期列看反、单位错 \\
\rowband
分季度 & 列示单季 = 累计相减 & 季度与累计期间混淆、追溯调整 \\
\bottomrule
\end{tabularx}
\end{table}

\subsection{异常信号}

程序按规则判出八类值得关注的信号（表~\ref{tab:fin-signals}），每条附带触发它的真值表单元格。信号只说“值得看”，不说“有问题”，
原因要由主笔会话回到年报附注和外部资料里去找。例如茅台 2025 年触发了“经营现金流显著低于净利润”（净现比 0.75，经营现金流同比下降 33.46\%），
主笔会话在研报中把它解释为经销商预收款减少、渠道打款意愿下降的先行指标。数字由程序给出，解释由模型完成，这就是“事实 → 推论”的分工。

\begin{table}[htbp]
\centering
\caption{八类异常信号及三家公司的触发情况}
\label{tab:fin-signals}
\small
\begin{tabularx}{\linewidth}{@{}c l L c c c@{}}
\toprule
\hd{编号} & \hd{类别} & \hd{含义} & \hd{茅台} & \hd{宁德} & \hd{招商} \\
\midrule
S1 & 非经常性损益 & 利润对非经常性损益依赖度较高 & & & 中 \\
\rowband
S2 & 利润质量 & 归母净利润增速显著高于扣非增速 & & & 中 \\
S3 & 利润与现金流背离 & 经营现金流显著低于净利润 & 中 & & \\
\rowband
S4 & 利润与现金流背离 & 净利润增长而经营现金流下降 & & & \\
S5 & 营运资本 & 应收账款增速显著快于收入 & & & \\
\rowband
S6 & 营运资本 & 存货增速显著快于收入 & & 高 & \\
S7 & 盈利能力 & 毛利率大幅变动 & & & \\
\rowband
S8 & 会计口径 & 存在会计政策或口径变化的披露 & 中 & 中 & 中 \\
\bottomrule
\end{tabularx}
\end{table}

\subsection{利润变动归因}

“归母净利润增长 18\%，主要受益于收入增长”是研报里最常见的句子，但“主要”往往没有依据。系统把两期利润的变化逐项拆成十个驱动：
收入规模效应、毛利率效应、销售/管理/研发/财务四项费用、其他经营损益、营业外收支、所得税、少数股东损益。各项加总严格等于归母净利润的变化，
闭合误差就是利润表恒等式的残差。这样正文里任何“某因素贡献了多少”的说法都能与之对账。以招商轮船为例，
归母净利润同比增长 17.71\%，扣非只增长 0.18\%，归因桥显示第一大正向驱动是“其他经营损益”（$+10.02$ 亿元），与 S1、S2 两条信号相互印证。

\subsection{实际 vs 预期、业绩预告与 EPS 质量}

提供一致预期和业绩预告文件时，程序计算实际值相对一致预期的差额与幅度、实际值落在预告区间的哪个位置，并做每股收益质量筛查：
检查是否存在使每股收益失真的项目，以及每股收益的超预期幅度是否明显大于收入、营业利润的超预期幅度（后者往往意味着利润来自非经营因素）。
程序只给幅度，不贴“超预期”“不及预期”的标签，措辞交给主笔会话结合市场语境决定。

\section{五家公司的结果}

表~\ref{tab:fin-five} 汇总五份真实文档的真值表规模与勾稽结果。勾稽共 68 项，67 项通过，1 项无法评估：
宇树科技的招股说明书没有分季度表，季度加总无从核对。没有一项失败。

\begin{table}[htbp]
\centering
\caption{五份真实文档的真值表与勾稽}
\label{tab:fin-five}
\small
\begin{tabularx}{\linewidth}{@{}l l r c L@{}}
\toprule
\hd{公司} & \hd{文档} & \hd{真值格} & \hd{勾稽} & \hd{程序判出的信号} \\
\midrule
贵州茅台 & 2025 年报 & 153 & 20/20 & 经营现金流显著低于净利润；会计政策变化 \\
\rowband
宁德时代 & 2025 年报 & 151 & 15/15 & 存货增速显著快于收入（高）；会计政策变化 \\
招商轮船 & 2025 年报 & 157 & 21/21 & 非经常性损益依赖；归母增速高于扣非；会计政策变化 \\
\rowband
中际旭创 & 2025 年报 & 134 & 8/8 & 存货增速快于收入；毛利率大幅变动；会计政策变化 \\
宇树科技 & 招股说明书 & 154 & 3/4 & 非经常性损益依赖（高）；毛利率大幅变动；会计政策变化 \\
\bottomrule
\end{tabularx}
\end{table}

\begin{story}[一个数字的完整来历：茅台 2025 年营业收入]
分析师会话在 \file{facts.json} 里写下：营业收入，2025 年，原文“168,838,102,514.79”，单位元，第 61 页，标签“其中：营业收入”，合并利润表。
证据校验在第 61 页按数字边界找到该串，接受。规则通道在第 6 页“主要会计数据”表第 1 行第 1 列取到同一个值，第 9 页和第 61 页作为佐证。
两通道一致，合并为 \code{rule+llm}，置信度高。真值表与 2024 年的 170,899,152,276.34 元比较，算出同比 $-1.2060\%$；
年报自己列示的同比是 $-1.21\%$，勾稽通过。估值模块以 168,838.10 百万元作为分部预测的基期。主笔在研报第一页写下
“全年实现营业收入 1,688.38 亿元，同比下降 1.21\%”，数字护栏在 $\halfulp=0.005$ 亿元的容差内找到对应值，放行。

这里还有一个口径判断：茅台的“营业总收入”比“营业收入”多出财务公司的利息收入约 32 亿元。分析师会话在备注里写明选“营业收入”的理由，
规则通道的别名表也优先取“营业收入”行。两条通道独立做出了相同的选择；如果不同，合并会报冲突，交回核对。
\end{story}
